\documentclass[12pt,a4paper]{article}
\usepackage[margin=1in]{geometry}

\begin{document}


\section{Analyzing Dispersion Characteristics of Fine Particulate Matter
in High-Density Urban Areas: A Study Using CFD Simulation
and Machine Learning}
Author: Daeun Lee, Caryl Anne M. Barquilla and Jeongwoo Lee (2025)\\

Draft version:\\

Terminology: \\
- PM2.5: Particulate Matter with a diameter of 2.5 micrometers (10$^{-6}$) or less. \\
- AR (Aspect Ratio): The ratio of the length to the width of an object (in this case, buildings). Note: in this case AR=distance between buildings/height of buildings\\
\\
Used Methods: \\
- XGBoost, Random Forest\\
- CFD: ENVI-met, OpenFOAM, and ANSYS Fluent(2023 R2) Note: Computational domain was constructed by AIJ guidelines \\
Results: \\
- 91-95\% accuracy in predicting PM2.5 concentrations using XGBoost. \\
Some notes:\\
Threshold effect of AR, where values exceeding 2.5 improve PM2.5 dispersion under high wind velocity. \\
In loop-road networks, centrally elevated buildings enhance ventilation, while in grid-road networks, taller perimeter buildings shield inner blocks from arterial emissions\\
Tested 3 Road Layout Configurations: 1 Grid road, 2 Loop road, 3 T-junction road\\
3 Block height configurations: 1 Enclosed height configuration, 2 Even height configuration, 3 Elevated core configuration\\
4 AR configurations: 1 AR=1, 2 AR=1.5, 3 AR=3, 4 AR=4.5\\
And also took into account the distance from arterial roads, side space between buildings, wind speed, atmospheric pressure\\

\textbf{Final version:}
This article authors investigated the dispersion of fine particulate matter
(PM2.5) in high-density urban environments using a combination of
computational fluid dynamics (CFD) simulations and machine learning models.
In this study, PM2.5 refers to particulate matter with a diameter of
2.5micrometers or less. The authors analyzed how urban morphology and
environmental conditions influence pollutant dispersion.

To model airflow and pollutant transport, several CFD platforms were used,
including ENVI-met, OpenFOAM, and ANSYS Fluent (2023 R2). The computational
domain was constructed according to the AIJ guidelines for pedestrian wind
environment simulations. In addition, machine learning models such as
XGBoost and Random Forest were applied to predict PM2.5 concentration levels
based on the simulation data.

The experimental setup considered multiple urban configurations, including
three road layout scenarios (grid, loop, and T-junction roads),
three building height configurations (enclosed height, even height, and
elevated core), and four aspect ratio (AR) scenarios ranging from AR=1 to
AR=4.5. In this study, the aspect ratio was defined as the distance between
buildings divided by the building height. Other environmental variables,
such as distance from arterial roads, side space between buildings, wind
speed, and atmospheric pressure, were also included in the analysis.

The results showed that the XGBoost model achieved high predictive
performance, reaching an accuracy of approximately 91-95\% in estimating
PM2.5 concentrations. Furthermore, the study identified a threshold effect
for the aspect ratio, where AR values greater than 2.5 significantly improve
urban ventilation. The findings also indicate that centrally elevated
buildings enhance ventilation in loop-road networks, while taller perimeter
buildings in grid-road layouts can help shield inner urban blocks from
arterial emissions.

\section{Air Pollution Dispersion Modelling in Urban Environment
Using CFD: A Systematic Review}

Author: Mariya Pantusheva, Radostin Mitkov, Petar O. Hristov and Dessislava Petrova-Antonova \\

Draft version:\\

Terminology: \\
- V\&V: Verification and Validation, methods used to check whether a model works correctly and whether its results can be trusted. \\
- RANS: Reynolds-Averaged Navier--Stokes equations, a simplified approach for simulating airflow where the average behavior of the flow is calculated. \\
- LES: Large Eddy Simulation, a more detailed method that simulates turbulence more accurately but requires much more computational power. \\
\\

Used Methods: \\
- A systematic literature review of about 90 scientific articles related to air pollution dispersion modelling in urban areas using CFD. \\
- The reviewed studies were analyzed using more than 190 modelling characteristics grouped into seven categories: general information, geometry, mesh, physics models, boundary conditions, solution setup, and verification/validation. \\
- The analysed papers were also compared with commonly used best-practice guidelines such as COST 732, AIJ guidelines \\
\\

Results: \\
- Review showed that the most frequently used CFD software in urban air pollution studies are ANSYS Fluent and OpenFOAM. \\
- Most studies used three-dimensional models, and the majority of simulations were performed using steady-state RANS models rather than LES models. \\
- For modelling pollutant transport, the passive scalar equation and the species transport model were the most commonly applied approaches. \\
\\

Some notes:\\
The authors found that around 12--34\% of the analysed papers did not report some important modelling details that are necessary to reproduce the study. \\
In addition, none of the reviewed studies performed verification and validation strictly according to the recommended standards. \\
Many papers focused mainly on methodological research rather than solving real applied urban air pollution problems. \\
Most of the studies analysed gaseous pollutants related to traffic emissions, while particulate matter and other pollution sources were investigated less frequently. \\
The review also highlights that modelling choices such as computational domain size, mesh resolution, turbulence model, and boundary conditions have a strong influence on the reliability of CFD simulation results. \\

\section{Application of CFD Modelling for Pollutant Dispersion at an
Urban Traffic Hotspot}

Author: Giannis Ioannidis, Chaofan Li, Paul Tremper, Till Riedel and Leonidas Ntziachristos \\

Draft version:\\

Terminology: \\
- ABL: Atmospheric Boundary Layer, the lowest part of the atmosphere that is directly affected by the ground surface. \\
\\

Used Methods: \\
- CFD modelling with OpenFOAM for simulating pollutant dispersion in a real urban traffic hotspot in Augsburg, Germany. \\
- Traffic emissions for NOx, CO, and PM were estimated using COPERT Street software based on traffic data from September 2018. \\
- Meteorological input data such as wind speed and wind direction were taken from the smartAQnet sensor network. \\
- The model used steady-state RANS, the standard k-$\varepsilon$ turbulence model, and the passive scalar transport equation for pollutant dispersion. \\
- A sensitivity analysis was carried out to select the most suitable computational mesh. \\
\\

Results: \\
- Developed CFD model produced high-resolution concentration maps for NOx, CO, and PM at street scale. \\
- Predicted values were in the same general range as observations and followed the daily traffic related pollution trends. \\
- Average daily deviation between modelled and measured values was about 11\% for NOx, 26\% for CO, and 27\% for PM. \\
- Reported regression coefficients were $r^2 = 0.82$ for NOx, $r^2 = 0.54$ for CO, and $r^2 = 0.91$ for PM. \\
\\

Some notes:\\
Study area was a highly trafficked part of Augsburg, Germany, with 3D urban model taken from OpenStreetMap. \\
Computational domain height was set above 6H and the upstream distance above 5H, following common CFD recommendations. \\
Model used seven road emission sources placed along the main traffic artery. \\
Medium computational mesh with 9 million cells was selected because it gave similar results to the fine mesh but required much less computational time. \\
Results showed that pollutant concentrations can vary strongly even within a small street area. In the analysed 20 $\times$ 60 m area, local deviations reached about +50\% and -18\% compared to the sensor location. \\
Study also showed that higher pollutant concentrations were associated with lower wind velocity, lower turbulent kinetic energy, and lower pressure areas inside the street canyon. \\



\section{Hazardous Atmospheric Dispersion in Urban Areas: a
Deep Learning Approach for Emergency Pollution
Forecast}

Author: Mouhcine Mendil, Sylvain Leirens, Patrick Armand, Christophe Duchenne

Draft version:\\

Terminology: \\

Used Methods: \\
- After training neural network can predict the dispersion of hazardous pollutants without running new CFD simulations. \\
\\

Results: \\
- Deep learning model was able to reproduce pollutant dispersion patterns generated by CFD simulations. \\
- The model significantly reduced computation time compared to traditional CFD simulations. \\
- Predictions of pollutant concentration fields were close to CFD results in most tested scenarios. \\
\\

Some notes:\\
Prooposed deep learning model acts as a surrogate model that approximates CFD results much faster. \\
Such models can help quickly estimate affected areas after accidental releases of hazardous substances. \\


\section{Hybrid Physics-Machine Learning Framework for Forecasting
Urban Air Circulation and Pollution in Mountain-Valley Cities}

Author: Lyazat Naizabayeva, Gulbakyt Sembina and Gulnara Tleuberdiyeva
Draft version:\\

Terminology: \\
- RMSE: Root Mean Square Error
- Urban boundary layer is lowest part of the atmosphere above the city where buildings, terrain and weather influence airflow. \\
- Boussinesq approximation: simplification of Navier-Stokes equations used in atmospheric flow simulations. \\
\\

Used Methods: \\
- Hybrid modelling approach combining Navier  stokes equations with machine learning methods. \\
- Machine learning module based on gradient boosting and decision tree methods for improving prediction accuracy. \\
- Model trained using meteorological observations and air quality data from 2024. \\
- The framework predicts wind speed, temperature and particulate matter concentrations for short-term forecasts (around 6-24 hours). \\
\\

Results: \\
- Hybrid model improved prediction accuracy compared to purely statistical models. \\
- Forecast accuracy increased by about 16\% in terms of $R^2$, and RMSE decreased by about 25\%. \\
- For 6 hour forecasts the model achieved RMSE $\approx$ 1.2 m/s for wind speed, 1.8°C for temperature, and 0.012 mg/m$^3$ for particulate matter. \\
\\

Some notes:\\
Increase of wind speed by 1 m/s reduced particulate matter concentration by about 10-15\% during winter nights. \\
Hybrid approach combines physical consistency of atmospheric equations with flexibility of machine learning models. \\
\end{document}